\documentclass[12pt, a0paper]{article}

% ===== 页面 =====
\usepackage[
    paperwidth=84.1cm,
    paperheight=118.9cm,
    margin=2cm,
    columnsep=2cm
]{geometry}

% ===== 中文 =====
\usepackage{ctex}

% ===== 双栏 =====
\usepackage{multicol}

% ===== 表格 =====
\usepackage[table]{xcolor}
\usepackage{colortbl}
\usepackage{multirow}
\usepackage{booktabs}

% ===== 其他 =====
\usepackage{tikz}
\usepackage{enumitem}
\usepackage{amsmath}
\usepackage{amssymb}

% ===== 颜色 =====
\definecolor{myblue}{RGB}{25,55,109}
\definecolor{myorange}{RGB}{200,100,20}
\definecolor{mygreen}{RGB}{30,120,70}
\definecolor{myred}{RGB}{180,40,40}
\definecolor{mygray}{RGB}{130,130,130}
\definecolor{mypurple}{RGB}{100,40,160}
\definecolor{lightblue}{RGB}{235,245,255}

% ===== 字体缩放 (A0) =====
\renewcommand{\normalsize}{\fontsize{22}{30}\selectfont}
\renewcommand{\large}{\fontsize{28}{36}\selectfont}
\renewcommand{\Large}{\fontsize{34}{42}\selectfont}
\renewcommand{\LARGE}{\fontsize{40}{48}\selectfont}
\renewcommand{\huge}{\fontsize{48}{56}\selectfont}
\renewcommand{\Huge}{\fontsize{56}{64}\selectfont}
\renewcommand{\small}{\fontsize{18}{24}\selectfont}
\renewcommand{\footnotesize}{\fontsize{15}{20}\selectfont}

% ===== 收紧间距 =====
\setlength{\parskip}{0.3em}
\setlength{\parindent}{0pt}
\setlength{\multicolsep}{0.5em}
\raggedbottom

% ===== 节标题 =====
\usepackage{titlesec}
\titleformat{\section}[hang]
    {\huge\bfseries\color{myblue}}
    {\thesection}
    {1em}
    {}
    [\vspace{-1mm}{\color{myblue}\rule{\columnwidth}{1.5pt}}]
\titlespacing{\section}{0pt}{0.5em}{0.3em}

\titleformat{\subsection}[hang]
    {\Large\bfseries\color{myblue!80}}
    {\thesubsection}
    {1em}
    {}
\titlespacing{\subsection}{0pt}{0.3em}{0.1em}

% ===== 列表 =====
\setlist{nosep, leftmargin=2.5em, itemsep=0.3em}

% ===== 占位图 =====
\newcommand{\placeholder}[2]{%
    \begin{center}
    \colorbox{mygray!12}{\parbox{0.9\columnwidth}{%
        \centering
        \vspace{1.5em}
        {\Large\bfseries\color{mygray} [ 占位图： #1 ]}
        \vspace{0.8em}

        {\small\color{mygray} #2}
        \vspace{1.5em}
    }}
    \end{center}
}

\begin{document}

% =====================
% 标题
% =====================
\begin{center}
    {\fontsize{58}{68}\selectfont\bfseries\color{myblue}
        EEG 解码模型的发展：\\[3mm]
        从 EEGNet 到 EEG Foundation Model 的复现与分析\\[8mm]
    }

    {\fontsize{30}{38}\selectfont\color{mygray}
        基于 EEGNet、EEG Conformer 和 CBraMod 的比较研究\\[10mm]
    }

    {\fontsize{28}{34}\selectfont
        第一作者 (SUSTech ID) \hspace{2em} 第二作者 (ID) \hspace{2em} 第三作者 (ID)\\[4mm]
    }

    {\fontsize{20}{28}\selectfont\color{mygray}
        南方科技大学生物医学工程系 \hspace{3em}
        BMEB215 机器学习与医学工程应用 · 2026 Spring
    }

    \vspace{6mm}
    {\color{myblue}\hrule height 2.5pt}
    \vspace{6mm}
\end{center}

\begin{multicols}{2}

% ================================================================
% 1. 研究背景
% ================================================================
\section{研究背景 (Background)}

\subsection{EEG 解码的重要应用}

脑机接口 (BCI) · 睡眠分期 (Sleep Staging) · 神经疾病辅助诊断 · 认知状态监测

\subsection{EEG 解码面临的挑战}

\begin{itemize}
    \item 数据标注成本高 —— 临床、睡眠等任务需要专家标注
    \item 被试间差异显著 —— 个体差异导致模型难以泛化
    \item 设备差异显著 —— 不同设备、导联、采样率带来分布偏移
    \item 模型泛化能力有限 —— 单任务训练，跨任务复用困难
\end{itemize}

\subsection{项目目标}

复现三个具有代表性的 EEG 解码模型，分析其设计思路、表征能力和泛化能力的演化：

\begin{center}
\Large\bfseries
EEGNet \quad $\longrightarrow$ \quad EEG Conformer \quad $\longrightarrow$ \quad CBraMod
\end{center}

% ================================================================
% 2. 模型复现
% ================================================================
\section{模型复现 (Methods)}
\subsection{模型结构比较}

\begin{center}
\begin{tabular}{lccc}
    \toprule
    \textbf{模型} & \textbf{CNN} & \textbf{Transformer} & \textbf{预训练} \\
    \midrule
    EEGNet        & $\checkmark$ & -- & -- \\
    EEG Conformer & $\checkmark$ & $\checkmark$ & -- \\
    CBraMod       & $\checkmark$ & $\checkmark$ & $\checkmark$ \\
    \bottomrule
\end{tabular}
\end{center}

\vspace{1em}

\textbf{主要发现：}从 EEGNet 到 EEG Conformer 再到 CBraMod，模型表现出更强的表征能力、更好的跨任务泛化能力，以及更复杂的模型结构。

\subsection{EEGNet · CNN 时代}

\textbf{核心思想：}利用 EEG 领域知识设计轻量化卷积网络。

\textbf{关键模块：}Temporal Convolution $\rightarrow$ Depthwise Convolution $\rightarrow$ Separable Convolution

\textbf{特点：}局部时空特征提取，参数量小，适合小样本 EEG 任务。

\vspace{1em}

\subsection{EEG Conformer · CNN + Transformer 时代}

\textbf{核心思想：}结合 CNN 局部特征提取与 Transformer 全局建模。

\textbf{关键模块：}Convolution Module + Multi-Head Self-Attention

\textbf{特点：}同时建模局部特征与全局依赖，增强长时程信息捕获能力。

\vspace{1em}

\subsection{CBraMod · Foundation Model 时代}

\textbf{核心思想：}基于大规模脑电数据预训练，学习通用 EEG 表征。

\textbf{关键模块：}EEG Tokenization $\rightarrow$ Criss-Cross Attention $\rightarrow$ Self-Supervised Pretraining

\textbf{特点：}学习通用脑电表征，支持跨任务迁移，提升泛化能力。

\placeholder{模型架构对比}{%
    EEGNet / EEG Conformer / CBraMod 结构示意图。\\
    展示三个模型的关键模块和架构差异。\quad (等待绘制)}

% ================================================================
% 3. 复现结果
% ================================================================
\section{复现结果 (Results)}

\subsection{模型性能比较}

\placeholder{模型性能对比}{%
    三个模型在多个数据集上的性能比较。\\
    指标：Accuracy / F1 Score / Cohen's Kappa。\\
    \small (等待队友复现结果)}

\subsection{训练效率比较}

\placeholder{训练效率对比}{%
    三个模型的参数量 / 训练时间 / 推理速度对比。\\
    \small (等待队友复现结果)}



% ================================================================
% 4. 分析与讨论
% ================================================================
\section{分析与讨论 (Discussion)}

\subsection{观察一：模型能力持续提升}

\begin{center}
\Large\bfseries
局部特征 \quad $\longrightarrow$ \quad 全局依赖 \quad $\longrightarrow$ \quad 通用表征
\end{center}

从 CNN 局部时空特征，到 Transformer 全局时序依赖，再到预训练通用表征——每一步都提升了模型的表征空间和泛化边界。

\subsection{观察二：Foundation Model 成为新的趋势}

CBraMod 通过大规模预训练，利用海量无标签 EEG 数据学习通用表示，支持多任务迁移。这改变了 EEG 解码的训练范式：不再为每个任务单独训练，而是预训练后再适配。

\subsection{观察三：可解释性成为新的挑战}

传统 EEG 分析依赖 Alpha、Beta、Theta、Delta 等频段特征，容易与神经科学知识对应。而 Foundation Model 的表征是 Embedding 高维向量，虽然泛化能力更强，但其生理意义更难解释。

\begin{center}
\colorbox{myorange!10}{\parbox{0.9\columnwidth}{\centering\bfseries
    泛化能力 $\uparrow$ \hspace{2em} 但 \hspace{2em} 解释难度 $\uparrow$
}}
\end{center}

随着 Foundation Model 的发展，如何理解模型学到的内容，正在成为新的研究问题。

\placeholder{可解释性挑战示意图}{%
    传统 EEG 分析 (Alpha/Beta/Theta + 可解释生理量)\\
    $\downarrow$\\
    Foundation Embeddings (高维隐空间 · 泛化强但难以理解)\\
    $\downarrow$\\
    抽象程度递增，解释难度增大。\quad (等待绘制)}

% ================================================================
% 5. 未来工作
% ================================================================
\section{未来工作 (Future Work)}

\subsection{NeuroGate：重新引入神经科学知识}

\textbf{目标：}

\begin{center}
\Large\bfseries
Foundation Model \quad + \quad Neuroscience Prior
\end{center}

\textbf{思路：}CBraMod Embedding + 神经科学先验知识（频带能量、时域统计等）$\rightarrow$ Gated Fusion $\rightarrow$ 分类预测。

\vspace{1em}

\textbf{研究问题：}是否能够在保持 Foundation Model 泛化能力的同时，提高模型的可解释性？

\placeholder{NeuroGate 结构示意图}{%
    CBraMod Backbone + 神经科学先验\\
    $\downarrow$ Gated Fusion\\
    增强表征 (泛化 + 可解释性)\\
    \small (等待绘制)}

\placeholder{NeuroGate 实验结果}{%
    LP / FT 下 NeuroGate vs Baseline 的对比结果。\\
    \small (等待队友复现结果)}

% ================================================================
% 6. 总结
% ================================================================
\section{总结 (Summary)}

\begin{center}
\colorbox{lightblue}{\parbox{0.95\columnwidth}{\centering
    {\LARGE\bfseries
        泛化能力持续提升\par}
    \vspace{1.5em}
    {\large
        EEGNet/EEG Conformer $\rightarrow$ CBraMod\par}
    \vspace{1em}
    {\large
        但可解释性问题逐渐凸显\par}
    \vspace{1.5em}
    {\large\bfseries
        未来方向：Foundation Model + Neuroscience Prior\\[2mm]
        $\rightarrow$ 强泛化 + 强解释\par}
}}
\end{center}

% ================================================================
% 参考文献
% ================================================================
\vspace{1em}
\section*{参考文献}

{\footnotesize
\begin{enumerate}[leftmargin=2em]
    \item Lawhern et al. \textbf{EEGNet: A Compact Convolutional Neural Network for EEG-based Brain-Computer Interfaces}. \textit{J. Neural Eng.}, 2018.
    \item Song et al. \textbf{EEG Conformer: Convolutional Transformer for EEG Decoding and Visualization}. \textit{IEEE TNSRE}, 2023.
    \item Wang et al. \textbf{CBraMod: A Criss-Cross Brain Foundation Model for EEG Decoding}. \textit{ICLR}, 2025.
    \item NeuroGate / NeuroKE 内部参考：neural-prior gated fusion experiments.
\end{enumerate}
}

\end{multicols}

% =====================
% 底部
% =====================
\vspace{4mm}
{\color{myblue}\hrule height 2pt}
\vspace{3mm}
\begin{center}
{\footnotesize\color{mygray}
    BMEB215 · 机器学习与医学工程应用 · 2026 Spring \hspace{3em}
    南方科技大学 · 生物医学工程系
}
\end{center}

\end{document}
